\chapter{Introdução}

No cenário atual de globalização, o acesso à informação legal transcende as fronteiras nacionais, permitindo que indivíduos e comunidades entendam e interajam com sistemas jurídicos estrangeiros. O acesso a documentos legais traduzidos é essencial para garantir que cidadãos, imigrantes, estudantes e profissionais do direito possam compreender seus direitos e deveres em diferentes jurisdições, sem depender exclusivamente de tradutores especializados. A tradução de documentos jurídicos, portanto, assume um papel crucial na promoção da acessibilidade à justiça e na construção de pontes de entendimento entre diferentes culturas e legislações.

No entanto, a tradução jurídica não é uma tarefa trivial, pois envolve nuances que vão além da simples conversão linguística. O sistema jurídico carrega consigo terminologias específicas que podem ter interpretações e implicações diferentes, dependendo do contexto, ou então palavras cujo significado difere do uso cotidiano. A tradução precisa não apenas transmitir o conteúdo literal, mas também o contexto e a intenção das disposições legais, o que é particularmente importante para evitar mal-entendidos que possam resultar em consequências jurídicas.

Diante desse contexto, surge a necessidade de desenvolver soluções que unam pesquisa em Processamento de Linguagem Natural (PLN) e engenharia de software para produzir ferramentas práticas que atendam aos requisitos de fidelidade terminológica e preservação documental do meio jurídico. Este trabalho propõe o desenvolvimento de um sistema de tradução automática baseado em PLN, com ênfase na construção de um protótipo funcional \textit{full-stack} para tradução de documentos jurídicos entre o português, o inglês e o espanhol. Além da fundamentação teórica sobre contextualização e tradução automática, o projeto entrega um produto mínimo viável (MVP) que processa documentos HTML reais, preserva a estrutura documental (marcadores, parágrafos e placeholders) e integra técnicas modernas de linguagem — incluindo modelos de linguagem de larga escala (LLMs) via API — e uma camada opcional de contextualização por recuperação (RAG leve).

\section{Problemática}

Apesar das inovações em tradução automática, a complexidade da linguagem judicial e as sutilezas culturais continuam a ser desafios substanciais. A questão central que este trabalho busca abordar é: \textbf{como a tecnologia de Processamento de Linguagem Natural (PLN), incluindo modelos de linguagem de larga escala, pode ser aplicada em um sistema prático para capturar e refletir as nuances culturais e legais ao traduzir documentos jurídicos, preservando sua integridade estrutural?} As implicações éticas e legais da tradução inadequada de termos críticos ressaltam a necessidade de um enfoque mais refinado e consciente na tradução automática de documentos jurídicos.

Além das nuances linguísticas, existe um desafio técnico frequentemente negligenciado: \textbf{a preservação da estrutura documental}. Documentos jurídicos frequentemente vêm em formatos com marcações (HTML) que carregam significado — cabeçalhos, parágrafos legalmente numerados, citações e placeholders. A simples perda, reordenação ou mesclagem desses elementos durante o processo de tradução pode comprometer a integridade do documento. Outro aspecto prático é a consistência terminológica: garantir que termos jurídicos críticos sejam traduzidos de forma consistente ao longo de um mesmo documento ou de documentos relacionados. Portanto, a problemática deste trabalho combina dois eixos: (a) captura de nuances culturais e jurídicas por meio de PLN; (b) projeto e avaliação de uma pipeline orientada à engenharia capaz de traduzir documentos HTML preservando sua organização e coerência terminológica.

\section{Justificativa}

A crescente quantidade de documentos jurídicos disponibilizados online, aliada à necessidade de acesso à informação em um contexto global, torna urgente o desenvolvimento de soluções que considerem tanto as especificidades culturais quanto os requisitos práticos de preservação documental. A tradução automática pode ampliar o acesso à informação legal, mas somente se implementada de maneira a minimizar riscos decorrentes de perdas estruturais ou de inconsistência terminológica. Assim, investigar e implementar um sistema que una teoria (estudos sobre contextualização em PLN) e prática (um protótipo fullstack reprodutível) é uma contribuição relevante para as áreas de tradução automática, direito e engenharia de software.

\section{Objetivo}

\subsection{Objetivo geral}

Desenvolver e avaliar um protótipo fullstack de tradução automática para documentos jurídicos, baseado em técnicas de Processamento de Linguagem Natural (PLN) e integração com modelos de linguagem de larga escala, priorizando a preservação da estrutura HTML e a consistência terminológica entre o português e o inglês.

\subsection{Objetivos específicos}

Considerando o desenvolvimento do trabalho e o objetivo geral apresentado, destacam-se os seguintes objetivos específicos:

\begin{itemize}
    \item Projetar e implementar uma pipeline fullstack (frontend + backend) capaz de processar documentos jurídicos em HTML, segmentá-los, traduzir trechos preservando marcadores e placeholders, e remontar o documento traduzido;
    \item Integrar o protótipo com serviços de geração de linguagem (LLMs) via API e, opcionalmente, com uma camada de recuperação de contexto (RAG leve) baseada em glossário e corpus;
    \item Documentar e comparar tentativas experimentais com abordagens alternativas (por exemplo, modelos seq2seq disponíveis via Hugging Face), descrevendo as limitações práticas observadas;
    \item Avaliar a eficácia do sistema por meio de métricas automáticas (BLEU, WER, PER, TER) e análises qualitativas, verificando o ganho de consistência terminológica e a preservação estrutural;
    \item Publicar os resultados e oferecer recomendações práticas para a aplicação de PLN em contextos jurídicos, destacando limitações e direções futuras.
\end{itemize}

\section{Metodologia}

A metodologia deste trabalho combina investigação bibliográfica com desenvolvimento experimental aplicado. Inicialmente, será realizada uma revisão da literatura sobre tradução automática, contextualização em PLN e aspectos específicos da linguagem jurídica. Em seguida, procede-se ao desenvolvimento do protótipo, composto por:

\begin{itemize}
    \item um frontend desenvolvido com Quasar Framework e Vue.js para upload, visualização e avaliação de documentos;
    \item um backend em Python com FastAPI responsável pela orquestração do pipeline: leitura e indexação de HTML, segmentação em nós, armazenamento em banco, integração com backends de tradução (serviços LLM via API) e montagem de variantes traduzidas;
    \item um componente de contextualização (RAG leve) que recupera trechos relevantes de um glossário/corpus para adaptar a tradução quando aplicável;
    \item serviços de exportação (HTML/TXT) que preservam a estrutura original do documento;
    \item um módulo de avaliação que calcula métricas automáticas e permite comparações com traduções humanas de referência.
\end{itemize}

Serão realizados experimentos comparativos entre estratégias (baseline sem contexto, adapted com contexto recuperado) e também será relatado o histórico de tentativas com modelos seq2seq open-source (Hugging Face), descrevendo problemas práticos observados (limite de tokens, fragmentação de parágrafos, dificuldades na integração de glossários). Os resultados quantitativos e qualitativos serão analisados para identificar ganhos e limitações da solução proposta.

\section{Organização}

O presente trabalho está organizado da seguinte forma:

Capítulo 2: Revisão da Literatura — trabalhos relacionados, conceitos fundamentais de Tradução de Máquina, contextualização em PLN e desafios jurídicos.

Capítulo 3: Metodologia — detalhamento da abordagem adotada, descrição da pipeline, das ferramentas e das métricas de avaliação.

Capítulo 4: Proposta e Implementação — arquitetura do sistema, componentes do protótipo (frontend, backend, RAG, telemetria) e decisões de projeto.

Capítulo 5: Resultados e Discussão — experimentos executados, métricas, comparação entre variantes, análise qualitativa e limitações.

Capítulo 6: Conclusão e Trabalhos Futuros — síntese das contribuições, implicações práticas e sugestões para a continuidade do trabalho.
